% Part: computability
% Chapter: recursive-functions
% Section: trees

\documentclass[../../../include/open-logic-section]{subfiles}

\begin{document}

\olfileid{cmp}{rec}{tre}
\olsection{వృక్షాలు}

క్రమాలను సంఖ్యలతో, వాటి ధర్మాలు మరియు వాటిపై క్రియలను వాటి సంకేతాలపై
ఆదిమ పునరావృత్త సంబంధాలు, ప్రమేయాలతో సూచించినట్లే, కొన్నిసార్లు
వృక్షాలను సహజ సంఖ్యలుగా సూచించడం ఉపయోగపడుతుంది. ఇందుకు క్రమాలను,
వాటి సంకేతాలను వాడతాం. ఒక వృక్షం ఒకే శీర్షం (బహుశా ఒక గుర్తుతో)
కావచ్చు; లేదా అనేక ఉపవృక్షాలకు కలిపిన ఒక శీర్షం (బహుశా ఒక గుర్తుతో)
కావచ్చు. ఆ శీర్షాన్ని వృక్షపు \emph{మూలం} అంటారు; దానికి కలిపిన
ఉపవృక్షాలను దాని \emph{తక్షణ ఉపవృక్షాలు} అంటారు.

వృక్షాలను $\tuple{k,d_1,\dots,d_k}$ అనే క్రమంగా పునరావృత్తంగా
సంకేతీకరిస్తాం. ఇక్కడ $k$ తక్షణ ఉపవృక్షాల సంఖ్య; $d_1$, \dots,
~$d_k$ వాటి సంకేతాలు. శీర్షాలకు గుర్తులు ఉంటే వాటిని తక్షణ
ఉపవృక్షాల తరువాత చేర్చవచ్చు. కాబట్టి గుర్తు~$l$ గల ఒకే శీర్షంతో
ఉన్న వృక్షపు సంకేతం $\tuple{0,l}$; గుర్తు~$l_1$ గల మూలాన్ని,
గుర్తులు $l_2$, $l_3$ గల రెండు ఏకశీర్ష వృక్షాలకు కలిపిన వృక్షపు
సంకేతం $\tuple{2,\tuple{0,l_2},\tuple{0,l_3},l_1}$.

\begin{prop}
  \ollabel{prop:subtreeseq}
  సంకేతం~$t$ గల వృక్షపు ఉపవృక్షాలన్నిటి సంకేతాలను మూలకాలుగా గల
  క్రమపు సంకేతాన్ని తిరిగి ఇచ్చే ప్రమేయం $\fn{SubtreeSeq}(t)$ ఆదిమ
  పునరావృత్త ప్రమేయం.
\end{prop}

\begin{proof}
  మొదట $\fn{ISubtrees}(t)=\fn{subseq}(t,1,(t)_0)$ ఆదిమ పునరావృత్త
  ప్రమేయమనీ, వృక్షం~$t$ యొక్క తక్షణ ఉపవృక్షాల సంకేతాలను అది తిరిగి
  ఇస్తుందనీ గమనించండి. ఇప్పుడు, వృక్షం~$t$లో మూలం నుంచి అత్యధికంగా
  $n$ స్థాయుల
  దూరంలో ఉన్న ఉపవృక్షాలన్నిటి క్రమాన్ని గణించే సహాయక ప్రమేయం
  $\fn{hSubtreeSeq}(t,n)$ను నిర్వచించవచ్చు. మూలం నుంచి అత్యధికంగా
  $0$ స్థాయుల దూరంలో ఉన్న ఉపవృక్షాల క్రమం---అంటే $t$ యొక్క మూలం
  వద్ద మొదలయ్యేది---కేవలం $t$తో కూడిన క్రమం. అత్యధికంగా $n+1$
  స్థాయుల దూరంలో ఉన్న వృక్షం~$t$ యొక్క ఉపవృక్షాల క్రమాన్ని పొందడానికి,
  ఇప్పటికే ఉన్న $n$ స్థాయి ఉపవృక్షాల క్రమానికి, వాటిలోని ప్రతి వృక్షపు
  తక్షణ ఉపవృక్షాలన్నిటి క్రమాన్ని అనుసంధానిస్తాం.
  \sourcecorrection{OLTECRREM-008}{ఆంగ్ల మూలం $\fn{hSubtreeSeq}(t,n)$ను మూలం నుంచి సరిగ్గా $n$ స్థాయుల దూరంలోని ఉపవృక్షాల క్రమంగా వర్ణించింది; కానీ దాని పునరావృత్త సమీకరణ పూర్వస్థాయిలను ప్రతిసారీ ఉంచుతుంది. సమీకరణకు, చివరి ``అన్ని ఉపవృక్షాలు'' నిర్ధారణకు అనుగుణంగా ``అత్యధికంగా $n$ స్థాయులు''గా సరిచేశాం.}
  ఇవన్నీ జాబితా చేయడానికి, $f(x)$ క్రమాల సంకేతాలను తిరిగి ఇచ్చే
  ఆదిమ పునరావృత్త ప్రమేయమైతే,
  $g_f(s,k)=f((s)_0)\concat\dots\concat f((s)_{k-1})$ కూడా ఆదిమ
  పునరావృత్త ప్రమేయమని గమనించండి:
    \begin{align*}
      g_f(s, 0) & = \emptyseq\\
      g_f(s, k+1) & = g_f(s, k) \concat f((s)_k)
    \end{align*}
    ఉదాహరణకు, $s$ వృక్షాల క్రమమైతే
    $h(s)=g_{\fn{ISubtrees}}(s,\len{s})$ అనేది $s$లోని మూలకాల
    తక్షణ ఉపవృక్షాల క్రమాన్ని ఇస్తుంది.
    \sourcecorrection{OLTECRREM-009}{ఆంగ్ల మూలం $g_f(s,k)$ను $0$ నుంచి $k$ వరకూ $k+1$ విలువల అనుసంధానంగా నిర్వచించి, తరువాత $g_{\fn{ISubtrees}}(s,\len{s})$ను పిలిచింది; దాంతో క్రమపు పరిధి దాటిన $(s)_{\len{s}}$ కూడా చేరుతుంది. మొదటి $k$ మూలకాలనే అనుసంధానించే ఖాళీ ఆధారం, $k$వ మూలకాన్ని చేర్చే ఉత్తరగామి దశను పునరుద్ధరించాం.}
    దీనితో $\fn{hSubtreeSeq}$ను ఇలా నిర్వచించవచ్చు:
    \begin{align*}
      \fn{hSubtreeSeq}(t, 0) & = \tuple{t} \\
      \fn{hSubtreeSeq}(t, n+1) & = \fn{hSubtreeSeq}(t, n) \concat
      h(\fn{hSubtreeSeq}(t, n)).
    \end{align*}
    సంకేతం~$t$ గల వృక్షంలోని ఉపవృక్షాల గరిష్ఠ స్థాయి, అంటే మూలానికి,
    ఒక ఆకు శీర్షానికి మధ్య గరిష్ఠ దూరం, సంకేతం~$t$చే అవధి చేయబడుతుంది.
    కాబట్టి సంకేతం~$t$ గల వృక్షపు ఉపవృక్షాలన్నిటి సంకేతాల క్రమం
    $\fn{hSubtreeSeq}(t,t)$తో లభిస్తుంది.
\end{proof}

\begin{prob}
  \olref[cmp][rec][tre]{prop:subtreeseq} నిరూపణలోని
  $\fn{hSubtreeSeq}$ నిర్వచనం సాధారణంగా పునరుక్తులను చేర్చుతుంది.
  ఫలిత జాబితాలో ప్రతి ఉపవృక్షపు సంకేతం ఒక్కసారి మాత్రమే వచ్చేలా చేసే
  ప్రత్యామ్నాయ నిర్వచనాన్ని ఇవ్వండి.
\end{prob}

\end{document}
